% Weekly report 1. Paste into your own document. Plain LaTeX -- no extra packages.

\section*{Project plan}

The plan below is preliminary and will be revised with the supervisors as the project
develops. Risk is graded 1 (routine) to 5 (may not succeed as specified). The numbered
objectives are those of the project brief.

\begin{center}
\begin{tabular}{clc}
\hline
Week & Activity & Risk \\
\hline
 1     & Project kick-off; framework and dataset validated end to end                    & 1 \\
 1--2  & Clinical background: anatomy, dominance, variants, CAD, CCTA (obj.\ 1--3)       & 1 \\
 2--4  & Literature review: segmentation and topological description (obj.\ 4)           & 2 \\
 2--3  & Data chapter: cohort composition and centerline validation                      & 1 \\
 3--5  & Reproduce the published segmentation benchmark from delivered weights (obj.\ 5) & 2 \\
 4--7  & Topological feature extraction: branching, angles, tortuosity (obj.\ 6--7)      & 3 \\
 6--9  & Cohort-wide feature extraction and population statistics (obj.\ 8)              & 3 \\
 9--13 & Outlier and extreme-value detection over the feature distribution (obj.\ 9)     & 4 \\
12--16 & Association with patient outcomes, \emph{if outcome data is obtained} (obj.\ 10)& 5 \\
12--18 & Hemodynamic surrogates as a fallback functional endpoint                        & 4 \\
16--20 & Consolidation, revision and thesis writing                                      & 2 \\
\hline
\end{tabular}
\end{center}

\section*{Progress this week}

\paragraph{Dataset validated.} The delivered dataset comprises 800 usable cases, each with a
CT volume, a multi-label segmentation, left and right centerlines and a lumen surface. The
segmentation labels 1--14 were resolved to artery names by cross-referencing the centerlines,
which carry both the integer label and the artery name, and verified against the voxel data.

\paragraph{Centerline topology validated cohort-wide.} All 1600 centerline sides were rebuilt
as rooted trees and checked against four invariants. All 1600 pass; 1584 are a single clean
rooted tree. The remaining sides are anatomical or annotation variants and are individually
accounted for rather than discarded.

\paragraph{Pretrained weights verified.} The published benchmark weights were staged and each
was loaded into the architecture its configuration declares, matching exactly. Eight of the
nine benchmark methods can therefore be reproduced without retraining. One cannot: the
ImageCAS three-stage baseline is missing its first-stage checkpoint.

\section*{Questions for the supervisors}

\begin{enumerate}
  \item \textbf{Outcome data (objective 10).} The only outcome-like field currently available
    is a binary disease flag. Is a richer outcome linkage available, or invasive fractional
    flow reserve for any subset? This determines whether objective 10 is feasible as stated
    or whether a computed hemodynamic endpoint should replace it. \emph{This is the
    highest-risk dependency in the plan.}
  \item \textbf{Missing Stage-1 checkpoint.} The delivered weights include four of the five
    stages of the ImageCAS baseline. Does the first-stage coarse checkpoint exist?
  \item \textbf{Annotation scope on the right side.} The proximal branches of the right
    coronary artery -- conus, sinoatrial nodal and acute marginal -- are not labeled anywhere
    in the dataset, whereas five left-side branches are. Was this deliberate? It affects
    whether left-versus-right branch counts can be compared.
  \item \textbf{Novelty check.} A literature search found no coronary-specific work applying
    outlier or extreme-value detection to topological features. Does that match the
    supervisors' knowledge of the field?
\end{enumerate}

\section*{Plan for next week}

Complete the literature review for objective 4; draft the data chapter; reproduce one
published benchmark row end to end from the delivered weights.
